% ECONOMICS_PROBLEM: {"id": "C14-91", "title": "Minimax rates for stochastic interventions beyond the parametric regime", "jel_code": "C14", "source_id": "OP-0525"}
% !TeX program = xelatex
\documentclass[11pt,a4paper]{article}
\usepackage[margin=23mm]{geometry}
\usepackage{fontspec}
\setmainfont{Latin Modern Roman}
\usepackage{amsmath,amssymb,mathtools}
\usepackage{xcolor,enumitem,booktabs,longtable,array,xurl,hyperref}
\definecolor{muted}{HTML}{53636C}
\hypersetup{colorlinks=true,urlcolor=blue,linkcolor=blue}
\setlength{\parindent}{0pt}
\setlength{\parskip}{5pt}
\newcommand{\R}{\mathbb R}
\newcommand{\E}{\mathbb E}
\newcommand{\N}{\mathbb N}
\newcommand{\ind}{\mathbf 1}
\newcommand{\Var}{\operatorname{Var}}
\newcommand{\Cov}{\operatorname{Cov}}
\newcommand{\supp}{\operatorname{supp}}
\DeclareMathOperator*{\argmin}{arg\,min}
\DeclareMathOperator*{\argmax}{arg\,max}
\newcommand{\entrymeta}[3]{{\sffamily\small\color{muted}\textbf{\texttt{#1}}\quad|\quad #2\par #3\par}}
\newcommand{\Problemref}[1]{\texttt{#1}}
\newcommand{\JELref}[2]{\texttt{#2} (JEL \texttt{#1})}
\begin{document}
\section*{Minimax rates for stochastic interventions beyond the parametric regime}
\textbf{Problem ID:} \texttt{C14-91}\quad \textbf{JEL:} \texttt{C14}\par
% BEGIN ECONOMICS STATEMENT
\label{problem:OP-0525}
\label{beyond:CI-46}
\entrymeta{CI-46}{Econometric theory and causal inference}{Research agenda formalization}
\subsection*{Definitions and assumptions}
Use $n$ independent binary-baseline observations with $Y\in\{0,1\}$, consistency, conditional exchangeability, and fixed overlap $\varepsilon\leq e\leq1-\varepsilon$. Fix $d,L,c,C$ and positive $s_e,s_m,s_p$. Assume $e\in\mathcal H_d^{s_e}(L)$, $m_0,m_1\in\mathcal H_d^{s_m}(L)$, and $p\in\mathcal H_d^{s_p}(L)$ with $c\leq p\leq C$. For a fixed $\delta\in(0,\infty)\setminus\{1\}$, an incremental propensity intervention assigns treatment with probability
\[
 q_\delta(x)=\frac{\delta e(x)}{1-e(x)+\delta e(x)}.
\]
The intervention uses a new randomization conditional on $X$, independent of the potential outcomes. Its identified mean is
\[
 \psi_\delta(P)=
 \int \frac{\delta e(x)m_1(x)+(1-e(x))m_0(x)}
 {1-e(x)+\delta e(x)}p(x)\,dx.
\]
\subsection*{Formal question}
For the displayed model $\mathcal P$ and squared-error loss, determine
\[
 R_{n,\delta}^\star
 =\inf_{\widehat\psi_\delta}\sup_{P\in\mathcal P}
 \E_P[(\widehat\psi_\delta-\psi_\delta(P))^2]
\]
as $n\to\infty$, with matching lower bounds and attainable upper bounds. Describe its dependence on $d,s_e,s_m,s_p$, including cases in which it is larger than order $n^{-1}$. Compare interventions defined through the unknown propensity with a supplied, fixed policy $q(x)$, whose target is $\int\{q m_1+(1-q)m_0\}p$; these are different statistical functionals.
\subsection*{Scope and status}
The intervention, loss, and nuisance classes make one research direction precise. The excluded case $\delta=1$ gives $\psi_1(P)=\E_P[Y]$, directly estimable by the sample mean, and therefore must not support a universal claim of slower rates. The source concerns optimality outside the root-$n$ regime.
\subsection*{References for the source of the problem}
Carlos Cinelli, Avi Feller, Guido Imbens, Edward Kennedy, Sara Magliacane, and Jose Zubizarreta, \emph{Challenges in Statistics: A Dozen Challenges in Causality and Causal Inference}, arXiv:2508.17099v1, 23 August 2025, Section 3.5.3, \emph{New parameters}.
\url{https://arxiv.org/html/2508.17099v1#S3.SS5.SSS3}

\subsection*{Common conventions: Source mathematical conventions}

\textbf{Well-defined objects.}
All probability spaces and measurable variables are specified or inherited
from a local statistical model. Expectations used as finite targets require
the stated integrability. A decision rule or estimator is measurable with
respect to its available information; randomized rules may use an auxiliary
independent random variable. Norms without a subscript are Euclidean in
finite-dimensional spaces. An optimizer is requested only when attainment
is assured; otherwise the question concerns the optimal value and
$\varepsilon$-optimal admissible rules for every $\varepsilon>0$.

\textbf{Identification.}
A structural model $m\in\mathcal M$ implies an observable law
$P_m^{\rm obs}$ and target $\theta(m)$. The target is point identified on
$\mathcal M$ if
\[
 P_{m_1}^{\rm obs}=P_{m_2}^{\rm obs}
 \quad\Longrightarrow\quad\theta(m_1)=\theta(m_2)
 \qquad(m_1,m_2\in\mathcal M).
\]
When this fails, the sharp identified set is
\[
 \Theta(P;\mathcal M)=\{\theta(m):m\in\mathcal M,
                                  P_m^{\rm obs}=P\}.
\]
Specifying an intervention or estimand does not establish identification.
Positivity, exclusion, causal sufficiency, measurement restrictions, and
transport assumptions are substantive local inputs, never automatic
consequences of notation.

\textbf{Minimax risk and nontrivial inference.}
For a statistical experiment with data law $P^{(n)}$, target $\theta(P)$,
and nonnegative loss $\ell$, define
\[
 \mathcal R_n(\mathcal P,\ell)
   =\inf_{\widehat\theta_n}\sup_{P\in\mathcal P}
      \E_{P^{(n)}}\ell(\widehat\theta_n,\theta(P)).
\]
The infimum is over measurable estimators. A sharp rate is a sequence
$r_n$ for which $0<c\leq C<\infty$, independent of $n$, satisfies
$c r_n\leq\mathcal R_n\leq C r_n$ for all sufficiently large $n$.
Squared-error parametric risk is of order $n^{-1}$; the corresponding
root-mean-squared error is of order $n^{-1/2}$. These different conventions
are distinguished locally. A uniform asymptotic confidence region $C_n$
must satisfy
\[
 \liminf_{n\to\infty}\inf_{P\in\mathcal P}
     \Pr_{P^{(n)}}\{\theta(P)\in C_n\}\geq 1-\alpha,
 \qquad 0<\alpha<1,
\]
together with an informative diameter or expected-width requirement.
Coverage alone permits the vacuous whole parameter space. Testing questions
similarly impose both size control and power against a declared separated
alternative class.

\textbf{Binary causal baseline.}
When an entry explicitly invokes this baseline, one observes independent
copies of $O=(X,W,Y)$, with $W\in\{0,1\}$ and potential outcomes $Y(0),Y(1)$.
Assume consistency $Y=Y(W)$, no interference, conditional exchangeability
$(Y(0),Y(1))\perp W\mid X$, and overlap
$\epsilon\leq e(X)\leq1-\epsilon$ almost surely for a fixed
$0<\epsilon<1/2$, where
\[
 e(x)=\Pr(W=1\mid X=x),\qquad
 m_w(x)=\E[Y\mid W=w,X=x],\qquad
 \tau(x)=m_1(x)-m_0(x).
\]
These assumptions identify conditional mean effects almost everywhere
under the covariate law. Pointwise or uniform effects require appropriate
regularity and a specified support. Continuous treatment, longitudinal
data, adaptive experiments, and network interference require their own
assumptions in place of this baseline. Bounded outcomes, densities, and
smoothness are additional assumptions only where stated.

\textbf{Function classes.}
On $[0,1]^d$, $\mathcal H_d^s(L)$ is the isotropic Hölder ball of radius
$L>0$: put $k=\lceil s\rceil-1$ and $\gamma=s-k\in(0,1]$, bound derivatives
through order $k$, and require the order-$k$ derivatives to be
$\gamma$-Hölder with constant at most $L$. Thus integer orders use a
Lipschitz final derivative convention. The class
$\operatorname{BV}_d(L)$ consists of functions $f\in L^1((0,1)^d)$ whose
distributional derivative is a finite vector measure and for which
$\|f\|_{L^1}+|Df|((0,1)^d)\leq L$.
An $L^\infty$ bound or a particular pointwise version is imposed separately
when needed. Sparse and neural-network classes require a declared
dictionary or architecture and coefficient bounds.

An anisotropic H\"older class with declared block smoothness indices means
coordinatewise H\"older regularity in each block, uniformly in the other
blocks, with all corresponding derivative and H\"older bounds at most
the stated radius. Derivative orders and the integer-order convention in
each block are those just given. Mixed-derivative bounds are additional
restrictions only when explicitly stated.

\textbf{Decision and computational questions.}
Policy regret compares admissible feasible policies under the same law and
constraints. In computational entries, finite data and thresholds have
rational encodings; running time is measured in their total bit length.
Real-valued equilibrium search uses the explicitly declared algebraic or
FIXP encoding. An approximation ratio for maximization is written
$\operatorname{OPT}/\operatorname{ALG}$ and for minimization as
$\operatorname{ALG}/\operatorname{OPT}$, with zero optima and infeasible
instances handled locally. Historical approximation bounds are not
automatically current bounds.
% END ECONOMICS STATEMENT
\end{document}
